\documentclass[10pt,a4paper]{article}
\usepackage[a4paper,left=1.45cm,right=1.45cm,top=1.0cm,bottom=0.9cm]{geometry}
\usepackage{newtxtext,newtxmath}
\usepackage{microtype}
\usepackage{xcolor}
\usepackage{enumitem}
\usepackage{titlesec}
\usepackage[hidelinks]{hyperref}

\definecolor{accent}{HTML}{8C1515}
\definecolor{tum}{HTML}{0065BD}
\definecolor{muted}{HTML}{505050}

\titleformat{\section}[runin]{\normalfont\bfseries\color{accent}}{}{0pt}{}[.]
\titlespacing*{\section}{0pt}{2.6pt}{0.45em}
\setlength{\parindent}{0pt}
\setlength{\parskip}{1.2pt}
\setlist{nosep,leftmargin=1.25em,itemsep=0.3pt,topsep=0.5pt}
\linespread{0.96}
\pagestyle{empty}

\newcommand{\hyp}[1]{\textbf{\color{tum}#1}}
\newcommand{\cit}[1]{\textsuperscript{\color{muted}[#1]}}
\newcommand{\step}[1]{\textbf{\color{tum}#1}}

\begin{document}

{\footnotesize\color{muted} Master's thesis exposé · Draft 18.10.2026 · Intended supervision: Prof.\ Grabmair, TUM Legal Tech \hfill Mert Ayvazoglu · M.Sc.\ Information Systems, TUM}\\[-2pt]
{\color{accent}\rule{\linewidth}{0.8pt}}\\[0pt]
{\fontsize{13}{15}\selectfont\bfseries Right Outcome, Wrong Reasons? Outcome versus Proof-Step Rewards\\[-1pt] for Faithful Statutory Reasoning in Language Models}\\[1.5pt]
{\small\itshape A controlled study of the RL training recipe on executable statutes, where a solver decides which facts are operative.}\\[-3pt]
{\color{accent}\rule{\linewidth}{0.4pt}}

\section{Motivation}
An LLM whose stated reasons are not the ones that produced its answer earns full reward under reinforcement learning with verifiable rewards (RLVR), as long as the answer is right. A German authority may not decide this way: it must state ``the essential factual and legal grounds'' that ``moved'' it to its decision (§\,39(1) sentence~2 VwVfG; tax authorities owe a weaker duty under §\,121(1) AO). The AI Act lets persons affected by decisions based on high-risk legal AI (Annex~III No.\,8(a)) obtain ``clear and meaningful explanations'' of the system's role and the decision's main elements (Art.\,86(1)). Such explanations help only if they are faithful, and LLMs applied to legal tasks fall short. Reasoning models collapse on integrated case-based legal analysis (11.46 to 33.99\% accuracy), and four of five spend more thinking tokens on wrong answers than on right ones\cit{1}. On contract entailment, models present answers as solver-derived that the solver, when run, does not return (15.3 to 52.5\% of cases, the upper value including invalid inputs)\cit{2}. Models name the right authority while their verdict does not follow it\cit{3}.

\section{Gap}
Outside law, RLVR added after SFT lowers reasoning faithfulness\cit{4} (its law task is multiple choice), and verified process supervision restores reasoning quality in chess and math\cit{5}. In law, executable statutes have served to evaluate LLMs\cit{6,7}, as an outcome reward for LLM-generated Prolog\cit{8} and as a proposed but untested reward\cit{9}. Legal RL rewards score outcomes or rubric- and structure-based steps\cit{10}; none checks the stated derivation against an executable statute. The closest design\cit{11} shapes GRPO with solver certificates in synthetic rule worlds with structured outputs; real statutes, free-text reasoning and transfer remain open. The thesis studies the \emph{training recipe} at 4B scale, where the reward is the variable.

\section{Research question}
\emph{Does outcome-only RL lower the causal faithfulness of LLM statutory reasoning, and do rewards for solver-verified proof steps prevent this at an acceptable accuracy cost?}
\begin{itemize}
  \item[\hyp{H1}] Relative to the base model and to answer-only SFT on the same items, outcome-GRPO raises accuracy and lowers counterfactual faithfulness (CF) on jointly correct items. A pilot tests this first.
  \item[\hyp{H2}] Proof-step GRPO yields higher CF than outcome-GRPO. Pre-registered as a trade-off: up to 3 points accuracy loss is acceptable.
  \item[\hyp{H3}] (secondary) The CF advantage of H2 persists on held-out rule families.
\end{itemize}

\section{Methodology}
\step{Environment.} Four rule families with executable ground truth and editable case facts: SARA income tax (§§\,1, 2, 63, 68, 151, 152, 7703) and SARA unemployment tax (§§\,3301, 3306)\cit{7}, OpenExempt bankruptcy exemptions\cit{6} and DeonticBench Airline baggage policies\cit{8}. A fact sampler is built for SARA's event-based Prolog cases; OpenExempt generates cases on demand. Templates render facts; an LLM paraphrase is kept only if parsing it back and re-running the solver reproduces the label. Variants with renamed sections and defined terms, in statute text and code, control for memorisation (SARA is public since 2020).
\step{Proof semantics.} A SWI-Prolog meta-interpreter records proof trees, including finitely failed exception branches under negation as failure. With several proofs, a node counts as \emph{necessary} only if it occurs in every proof; an arithmetic built-in counts as one step. OpenExempt's step targets are its solver's intermediate results (e.g.\ jurisdiction under the 730-day rule). The statutes are edited to be unambiguous: the study measures faithfulness to one encoding, not open texture.
\step{Arms} (Qwen3-4B in thinking mode, LoRA). Base; answer-only SFT; outcome-GRPO; SFT followed by outcome-GRPO, reported both to help\cit{12} and to hurt\cit{4} faithfulness; proof-step GRPO. The step reward credits a cited step only if it re-executes in the solver \emph{and} belongs to the gold tree, scored as precision and recall over its nodes, so padding with valid but irrelevant steps earns nothing. Outcome-GRPO already assigns implicit Monte-Carlo step credit\cit{13}, so the contrast isolates explicit solver verification. Translate-and-execute is an inference-only reference: executable law needs no prose reasoning, but most law is not formalised, so the executable families serve as a laboratory with known derivations.
\step{Metrics.} \emph{Primary, pre-registered:} CF, the balanced accuracy of ``the answer changes if and only if the solver's answer changes'' over solver-verified minimal pairs (operative versus irrelevant fact edits). Computed on items that both compared models answer correctly before the edit, within difficulty bins. \emph{Secondary, difficulty-matched:} CoT truncation and mistake insertion\cit{14} and token spend on wrong versus right answers. \emph{Reason adequacy} (necessary proof nodes stated in the explanation) is scored only on held-out families with tags stripped and tag-free prompting, so it does not re-measure the reward; two raters with US-law training validate it on 100 items stratified by family (target Krippendorff's $\alpha \geq 0.67$). Three seeds in-distribution, two per transfer fold; item-level bootstrap; minimum effect of interest 5 CF points.

\section{Pilot and plan}
\step{H1 pilot (weeks 1 to 2).} OpenExempt, 4,000 training and 400 test items; Qwen3-1.7B/4B; base, answer-only SFT, outcome-GRPO; 2 seeds; 25 to 35 H100-hours. H1 is supported if, against base, outcome-GRPO gains $\geq$5 accuracy points and loses $\geq$5 CF points on jointly correct items in both seeds (bootstrap 95\% CI excluding 0). A 2 to 5 point drop triggers a third seed. If CF rises or moves under 2 points, the thesis reports that and pursues H2/H3 as sample efficiency and transfer.
\step{Schedule.} M1: pilot, SARA sampler, meta-interpreter, pair generators; M2: Airline wrapper, main runs (3 RL arms $\times$ 3 seeds); M3 to M4: transfer (4 folds $\times$ 2 RL arms $\times$ 2 seeds), secondary metrics, raters; M5a: buffer; M5b to M6: writing. SARA folds keep the shared definitions (§§\,152, 7703) in training, so they test new top-level rules. Compute: about 350 H100-hours at about 10 per 4B run (LRZ to be secured; rented: USD 900 to 1,100); if SARA's long prompts cost more, transfer folds drop to one seed.
\step{Risks.} Unstable GRPO: expert iteration on proof-filtered samples. A competing paper on proof-step RL over real statutes: the CF measurement and transfer study remain.

\section{Contribution}
(i)~To my knowledge, the first controlled test on executable statutes of whether outcome-only RL lowers the faithfulness of legal reasoning, and whether proof-step rewards prevent it. (ii)~An open environment with solver-verified counterfactual pairs for four rule families. (iii)~CF and reason adequacy as proposed audit proxies for whether a legal AI's stated grounds are the operative ones. The thesis answers the call in [1] for ``post-training methodologies designed for legal reasoning'' and complements the proposal by Kant, Nabi et al.\ (CodeX) to train Prolog translators with interpreter feedback (arXiv:2502.17638).

\vspace{0pt}{\color{accent}\rule{\linewidth}{0.3pt}}\\[-1.5pt]
{\fontsize{6.6}{7.4}\selectfont\color{muted}
[1]~Zhang, Grabmair, Gray, Ashley: Thinking Longer, Not Always Smarter, CSLAW~'26.
[2]~Wang et al.: Know Your Limits, arXiv:2606.16118.
[3]~Sadhu, Arora, Seth: Cited but Not Consulted, arXiv:2610.12361.
[4]~Han, Lee, Do: RFEval, ICLR 2026.
[5]~Kim et al.: Correct Answers from Sound Reasoning (VPS), COLM 2026.
[6]~Servantez et al.: OpenExempt, Findings of ACL 2026.
[7]~Holzenberger, Blair-Stanek, Van Durme: SARA, NLLP 2020.
[8]~Dou et al.: DeonticBench, arXiv:2604.04443.
[9]~Heydari, Leowald: Closing the Loop, arXiv:2606.23913.
[10]~LegalΔ, JPO, LexReward: arXiv:2508.12281, 2608.29616, 2609.39071.
[11]~X.\ Zhang: NormWorlds-CF, arXiv:2607.03957.
[12]~Yu et al.: Outcome Rewards Do Not Guarantee Verifiable\dots\ Reasoning, ICML 2026.
[13]~Sullivan, Koller: GRPO is Secretly a Process Reward Model, ICML 2026.
[14]~Lanham et al., arXiv:2307.13702.\par}

\end{document}
